% !TEX encoding = UTF-8
% !TEX TS-program = pdflatex
% !TEX root = ../tesi.tex

%**************************************************************
\chapter{Validation}
\label{cap:chapter4}

\textcolor{magenta}{(I think it is a little bit unbalanced with respect to the other chapters. I would include it into the previous chapter. What do you think?}

The validation of the model behavior is a critical process to distinguish between plausible but incorrect results and valid ones.
In this work, validation is performed by comparing the simulated performance metrics with well-established theoretical results available in the literature. In particular, the\gls{BER} and\gls{SER} obtained through Monte Carlo simulations are evaluated against analytical expressions for M-\gls{PSK} modulation schemes over\gls{AWGN} channels, enabling the verification of both the quantitative accuracy of the results and the qualitative behavior of the system.
Through this validation framework, the model is assessed not only in terms of performance evaluation but also in terms of its consistency with established communication theory, providing confidence in the results presented in \chaptername~\ref{cap:chapter5}.
%**************************************************************

%**************************************************************
\section{Bit Error Rate}
\label{sec:ber}

\gls{BER} is one of the most widely used performance metrics for evaluating digital communication systems, defined as the ratio between the number of incorrectly received bits and the total number of transmitted bits.
\begin{equation}
    \label{eq:ber_definition}
    \text{BER} = \frac{N_{\text{errors}}}{N_{\text{total}}}
\end{equation}
\gls{BER} provides a measure of reliability of a communication link, capturing the cumulative impact of channel impairments such as noise, fading, interference, and hardware faults. In this work, it is evaluated as a function of the energy per bit to noise power spectral density ratio: $E_b/N_0$, to prefer a normalized metric that is independent of modulation order and bandwidth for fairer comparisons across different system configurations.

\section{Symbol Error Rate}

\gls{SER} is a performance metric defined as the probability that a transmitted symbol is incorrectly detected at the receiver. Unlike\gls{BER}, which operates at the bit level,\gls{SER} evaluates errors at the symbol level, making it useful for systems that employ multi-bit modulation schemes such as M-\gls{PSK}.

\begin{equation}
    \label{eq:ser_definition}
    \text{SER} = \frac{N_{\text{error symbols}}}{N_{\text{total symbols}}}
\end{equation}
As we have already specified,\gls{SER} is strongly dependent on the modulation order and the geometric arrangement of constellation points; as the constellation size increases, the symbols become more densely packed in the signal space, making them susceptible to noise and interference. Consequently, higher-order modulation schemes require higher energy-per-bit to maintain low\gls{SER}.

\section{Gray Coding}
The mapping between binary sequences and constellation symbols is crucial for determining the overall bit error performance. The most likely symbol errors occur between nearest neighbors in the signal constellation; a commonly adopted strategy to ensure that these neighbors differ by a single bit is Gray coding.
Under the assumption of gray-coded constellations and independent bit errors, the relationship between the\gls{SER} and the\gls{BER} can be approximated as:

\begin{equation}
    \label{eq:ber_ser_relation}
    \text{BER} \approx \frac{\text{SER}}{\log_2 M}
\end{equation}
This relationship becomes more accurate the higher the\gls{SNR} ratio is, where errors are dominated by nearest-neighbor symbol confusions.\\
In this work, gray mapping is utilized for all considered modulation schemes, ensuring consistent and fair comparison of\gls{BER} performance across different system configurations.
\begin{figure}[htbp]
    \centering
    \begin{subfigure}{0.45\textwidth}
        \includegraphics[width=\linewidth]{images/chapter4/2psk_gray_constellation.pdf}
        \caption{2-PSK gray coded constellation}
        \label{fig:2psk_gray_constellation}
    \end{subfigure}
    \hfill
    \begin{subfigure}{0.45\textwidth}
        \includegraphics[width=\linewidth]{images/chapter4/4psk_gray_constellation.pdf}
        \caption{4-PSK gray coded constellation}
        \label{fig:4psk_gray_constellation}
    \end{subfigure}
    \vspace{0.5cm}
    \begin{subfigure}{0.45\textwidth}
        \includegraphics[width=\linewidth]{images/chapter4/8psk_gray_constellation.pdf}
        \caption{8-PSK gray coded constellation}
        \label{fig:8psk_gray_constellation}
    \end{subfigure}
    \hfill
    \begin{subfigure}{0.45\textwidth}
        \includegraphics[width=\linewidth]{images/chapter4/16psk_gray_constellation.pdf}
        \caption{16-PSK gray coded constellation}
        \label{fig:16psk_gray_constellation}
    \end{subfigure}
    \caption{M-PSK (2, 4, 8, 16) gray coded constellations}
    \label{fig:mpsk_gray_constellation}
\end{figure}


\section{Theoretical M-PSK performance}
\label{sec:theoretical_mpsk}

For $M$-\gls{PSK} modulation over an\gls{AWGN} channel, the theoretical symbol error rate is:
\begin{equation}
    \label{eq:ser_mpsk}
    \mathrm{SER}(E_s/N_0) \approx 2\,Q\!\left(\sqrt{2\,\frac{E_s}{N_0}}\,\sin\frac{\pi}{M}\right), \quad M \geq 4
\end{equation}
where $Q(x) = \frac{1}{2}\,\mathrm{erfc}\!\left(x/\sqrt{2}\right)$ is the Q-function. For $M = 2$ (B\gls{PSK}), the expression is:
\begin{equation}
    \label{eq:ber_bpsk}
    \mathrm{BER}(E_b/N_0) = Q\!\left(\sqrt{2\,\frac{E_b}{N_0}}\right)
\end{equation}
Combining \equationname~\eqref{eq:ser_mpsk} with the Gray-coded\gls{BER} approximation of \equationname~\eqref{eq:ber_ser_relation} and substituting $E_s/N_0 = \log_2\!M \cdot E_b/N_0$ yields the theoretical\gls{BER} as a function of $E_b/N_0$ for $M \geq 4$:
\begin{equation}
    \label{eq:ber_mpsk}
    \mathrm{BER}(E_b/N_0) \approx \frac{2}{\log_2 M}\,Q\!\left(\sqrt{2\,\log_2\!M\cdot\frac{E_b}{N_0}}\,\sin\frac{\pi}{M}\right)
\end{equation}
These expressions serve as the analytical reference against which the Monte Carlo simulation results are compared in \chaptername~\ref{cap:chapter5}.

\section{M-PSK performance under a static phase offset}
\label{sec:ber_phase_offset}

The eavesdropper-directed reward terms of \chaptername~\ref{cap:chapter3} require the\gls{BER} experienced by a receiver whose decision regions are rotated by a static phase offset $\theta_e$ relative to the legitimate receiver, since the array response generally has a different phase at an eavesdropper direction $\varphi_e$ than at the target direction $\varphi^\star$. This section extends \equationname~\eqref{eq:ser_mpsk}--\eqref{eq:ber_mpsk} of \sectionname~\ref{sec:theoretical_mpsk} to this case.

Because the $M$-\gls{PSK} constellation is invariant under rotations of $2\pi/M$, and the decision regions are symmetric about each constellation point, any phase offset $\theta_e$ can be folded into an equivalent offset
\begin{equation}
    \delta = \min\!\left(\theta_e \bmod \frac{2\pi}{M},\;\; \frac{2\pi}{M} - \left(\theta_e \bmod \frac{2\pi}{M}\right)\right) \in \left[0,\, \frac{\pi}{M}\right]
    \label{eq:phase_offset_fold}
\end{equation}
measuring the distance of the rotated symbol from its nearest decision boundary (located at $\pm\pi/M$ from the symbol's nominal phase).

For $M = 2$, the classical result for a phase synchronization error applies directly: the signal's projection onto the decision axis is scaled by $\cos\delta$, giving
\begin{equation}
    \mathrm{BER}(E_b/N_0, \theta_e) = Q\!\left(\sqrt{2\,\frac{E_b}{N_0}}\,\cos\delta\right).
    \label{eq:ber_bpsk_phase_offset}
\end{equation}

For $M \geq 4$, the offset symbol sits closer to one of its two neighbouring boundaries than the other, so the two-sided nearest-neighbour approximation of \equationname~\eqref{eq:ser_mpsk} splits into two distinct terms, one per boundary:
\begin{align}
    \mathrm{SER}(E_s/N_0, \theta_e) \approx{} & Q\!\left(\sqrt{2\,\frac{E_s}{N_0}}\,\sin\!\Big(\max\big(\tfrac{\pi}{M} - \delta,\, 0\big)\Big)\right) \notag \\
    & {}+ Q\!\left(\sqrt{2\,\frac{E_s}{N_0}}\,\sin\!\left(\frac{\pi}{M} + \delta\right)\right)
    \label{eq:ser_mpsk_phase_offset}
\end{align}
where clamping the first argument at $0$ once $\delta \geq \pi/M$ (i.e.\ once the offset symbol has crossed the near boundary) saturates that term at its maximal value $Q(0) = 0.5$ rather than leaving it undefined. Combining \equationname~\eqref{eq:ser_mpsk_phase_offset} with the Gray-coded approximation of \equationname~\eqref{eq:ber_ser_relation} gives the\gls{BER} used for the eavesdropper reward terms of \chaptername~\ref{cap:chapter3}:
\begin{align}
    \mathrm{BER}(E_b/N_0, \theta_e) \approx{} & \frac{1}{\log_2 M}\Bigg[Q\!\left(\sqrt{2\log_2\!M\cdot\frac{E_b}{N_0}}\,\sin\!\Big(\max\big(\tfrac{\pi}{M}-\delta,\,0\big)\Big)\right) \notag \\
    & {}+ Q\!\left(\sqrt{2\log_2\!M\cdot\frac{E_b}{N_0}}\,\sin\!\left(\frac{\pi}{M}+\delta\right)\right)\Bigg]
    \label{eq:ber_mpsk_phase_offset}
\end{align}
clipped to $[0, 0.5]$ to remain a valid probability. As a sanity check, setting $\theta_e = 0$ (hence $\delta = 0$) collapses \equationname~\eqref{eq:ser_mpsk_phase_offset} to $2\,Q\!\left(\sqrt{2E_s/N_0}\sin(\pi/M)\right)$, recovering exactly \equationname~\eqref{eq:ser_mpsk}.

\section{Monte Carlo simulation}

The performance of the system is evaluated through Monte Carlo simulations, providing a flexible and widely adopted framework for estimating error rates. In this approach, the\gls{BER} and\gls{SER} are empirically estimated as:
\begin{equation}
    \widehat{\text{BER}} = \frac{N_{\text{bit errors}}}{N_{\text{total bits}}}, \quad
    \widehat{\text{SER}} = \frac{N_{\text{symbol errors}}}{N_{\text{total symbols}}}
\end{equation}
To characterize system performance,\gls{BER} and\gls{SER} are evaluated as functions of the energy-per-bit to noise power spectral density ratio, $E_b/N_0$, by varying the noise variance while keeping the transmitted signal energy constant. Given the stochastic nature of the simulation, the accuracy of the estimated error rates depends on the number of transmitted bits or until a predefined number of errors is observed, ensuring statistical reliability of the results.